\documentclass{ieeeaccess}

\usepackage[utf8]{inputenc}
\usepackage{cite}
\usepackage{amsmath,amssymb,amsfonts}
\usepackage{graphicx}
\usepackage{textcomp}
\usepackage{booktabs}
\usepackage{url}

\begin{document}

\history{Date of publication April 6, 2026.}
\doi{10.1109/ACCESS.XXXX.XXXXXXX}

\title{Anti-Sandbox Fingerprinting: Realistic Windows 11 VM Personalization for Dynamic Analysis}

\author{\uppercase{Raghottam}\authorrefmark{1},
\uppercase{Sumukha}\authorrefmark{1},
\uppercase{Snehal}\authorrefmark{1}}

\address[1]{Department of Computer Science and Engineering, RV College of Engineering, Bengaluru, India}

\corresp{Corresponding author: Raghottam}

\begin{abstract}
Automated malware analysis environments, commonly known as sandboxes, find themselves locked in a persistent, continuously evolving arms race with environment-aware malware. Today's sophisticated threat actors frequently execute advanced "anti-VM" or "anti-sandbox" checks, systematically scanning the host environment for signs of a sterile, freshly installed operating system that inexplicably lacks genuine everyday human activity. When such a synthetic environment is detected, modern malware intelligently alters its behavior, sleeps out the analysis timer, or ceases execution entirely, rendering the dynamic analysis pipeline useless. This paper presents ARC (Artifact Reality Composer), an extensively engineered, Python-based framework inherently designed to deeply personalize and modify Windows 11 virtual machine images at enterprise scale. By mathematically synthesizing highly coherent, profile-driven artifacts—including wide-ranging filesystem distributions, deep registry perturbations, realistic SQLite browser histories, and complex binary event log chains—ARC forcefully transforms a generic, empty virtual machine into a highly plausible "lived-in" system. Our foundational methodology emphasizes rigorous internal consistency across both spatial and temporal domains, ensuring that generated timestamps, dynamic user directories, and underlying hardware metadata logically align with specific, algorithmically defined usage personas (e.g., an "Enterprise Office User" or a "Software Developer"). Our empirical evaluation against multiple baselines demonstrates a massive 400\% to 10000\% increase in relevant artifact density across heavily targeted operating system subsystems. Crucially, ARC yields a mathematically significant improvement in the system's resistance to common sandbox detection heuristics, achieving over 98.4\% internal coherence. This research successfully blurs the lines between isolated, fragile analysis environments and genuine, actively utilized production machines, forcing continuous adaptation from adversarial evasion logic.
\end{abstract}

\begin{keywords}
Malware analysis, sandbox evasion, virtualization, dynamic analysis, artifact generation, Windows forensics
\end{keywords}

\maketitle

\section{Introduction}

Dynamic analysis remains a paramount, critical foundation of modern malicious software (malware) research, automated incident response, and active threat intelligence gathering. By safely detonating and observing unknown, potentially malicious code in a highly controlled virtualized environment, security researchers and automated pipelines can rapidly extrapolate network indicators of compromise (IoCs), critical file signatures, and deep payload execution paths. However, the operational efficacy of this dynamic approach is constantly undermined by advanced evasion techniques deployed by forward-thinking malware authors.

\subsection{The Sterile Sandbox Problem}
Modern malware variants routinely incorporate strict, multi-layered environmental checks into their primary execution payloads to rapidly determine if they are executing within a sandbox or an automated analysis pipeline. These deterministic checks range from simple user-mode API calls querying for hardware strings inherently associated with mainstream hypervisors (e.g., VirtualBox, VMware, QEMU, KVM) to incredibly advanced heuristic scans searching for a distinct lack of long-term user-specific data. A "sterile" sandbox—one deployed automatically from a pristine, clean OS image snapshot for momentary detonation—contains glaring, obvious systemic anomalies. It consistently features an entirely empty Recent Documents folder, completely untouched browser caches, identical logical disk partition serial numbers across millions of executions, and a severe, mathematically improbable lack of third-party software installation footprints. 

When these tell-tale signs of a computational synthetic environment are algorithmically identified, evasive malware often stalls its primary execution loop, sleeps for prolonged, un-analyzable periods, or cleanly exits the process tree entirely. In numerous observed APT campaigns, the malware may actively display benign "decoy" behavior designed explicitly to fool the automated tracking systems into classifying the payload as a clean, harmless executable.

\subsection{Limitations of Current Reactive Approaches}
The prevailing defensive methodology for counteracting these advanced environmental checks heavily relies upon human security analysts manually interacting with the system instance—creating a handful of dummy text documents, hastily browsing a few randomized web URLs, and altering superficial system settings before taking a master golden snapshot. While this approach succeeds in creating a somewhat believable base environment for a single, targeted execution, this strategy critically lacks the necessary scalability, determinism, and reproducibility required for continuous, automated, high-throughput threat intelligence systems. 

Furthermore, human interaction is strictly temporally constrained; generating the complex illusion of a workstation that has been naturally utilized over a span of several months is manually prohibitive and logically impossible to scale. Randomized, automated scripts simply generating files suffer from profound temporal clustering flaws, wherein thousands of files share exact second-precision creation timestamps, serving as an immediate algorithmic indicator of synthetic generation for anti-analysis routines.

\subsection{Scope and Contributions}
To comprehensively address these fundamental architectural limitations, we formally outline and rigorously evaluate the Artifact Reality Composer (ARC). ARC is an automated "VM Personalization" framework. Rather than acting as a traditional hypervisor-level cloaking mechanism (which is frequently defeated via CPU timing attacks), ARC operates securely and deeply on the raw data layer of a mounted Windows 11 Virtual Hard Disk (VHD).

This paper offers the following principal contributions to the field of malware analysis and digital forensics:
\begin{itemize}
\item We propose a fully decoupled, highly extensible service-oriented architecture capable of absolute deterministic "wear-and-tear" artifact generation without ever requiring the target VM to be actively booted.
\item We present a novel formal framework for generating internally coherent Spatio-Temporal identities. Our engine algorithmically aligns historical event distribution dates, filesystem usernames, and stochastically generated physical hardware characteristics.
\item We actively detail deep, offline Windows registry manipulation techniques designed to systematically isolate, scrub, and rewrite virtual machine indicators and seamlessly normalize hardware identification signatures.
\item We exhaustively evaluate our applied approach using a custom evaluation suite explicitly targeting Artifact Density, Internal Consistency Algorithms, and Sandbox Signal Detection Resistance, empirically demonstrating highly successful evasion of modern environment-probing heuristics.
\end{itemize}

\section{Background and Threat Model}

To accurately contextualize the necessity of environment personalization, it is crucial to understand the baseline capabilities of the adversary and the underlying architecture of modern dynamic analysis.

\subsection{The Baseline Malware Analysis Pipeline}
In a traditional sandbox environment, a virtual machine is instantiated from a golden, known-good snapshot. A submission mechanism places the malicious executable into the guest OS. Concurrently, a monitoring agent is injected. Over a pre-defined time window, the execution is monitored. After the timeout, the VM is forcefully terminated and reverted to restoring the golden snapshot for the next sample.

\subsection{Adversary Capabilities (The Threat Model)}
We operate under a threat model where the adversary possesses 'Environment Awareness'. The adversary can:
\begin{itemize}
\item \textbf{T1: Query Hardware Topologies.} Access WMI to pull SMBIOS strings, CPU attributes, and MAC address assignments.
\item \textbf{T2: Perform Artifact Enumeration.} Traverse the host filesystem and registry to quantify installed applications, active network profiles, and historical telemetry databases.
\item \textbf{T3: Execute Temporal Validation.} Interrogate MAC timestamps on user-generated data to identify unnatural temporal clustering.
\item \textbf{T4: Assess Behavioral Entropy.} Analyze the stochastic randomness of browser cache states and logical sequence event logs to detect algorithmic automation.
\end{itemize}

Our goal is active defeat of \textbf{T2, T3, and T4}—the core data-layer heuristics.

\section{Literature Review}

The ongoing, highly competitive arms race between malicious payload engineers and security analysts has witnessed a dramatic shift in recent years, pivoting aggressively from basic hypervisor detection towards sophisticated computational human-behavior analysis, a trend heavily and meticulously documented in contemporary systems security literature.

\subsection{Modern Evasion Taxonomies (2021-Present)}
Zhang et al. \cite{b1} authored an expansive, mathematically rigorous survey outlining the comprehensive taxonomy of dynamic analysis evasion. Their deep-dive research emphasizes definitively that malware fingerprinting of sandbox environments is no longer an isolated edge case restricted to advanced persistent threats (APTs). Instead, it operates as standard operational procedure across almost all commodity malware families spanning highly commoditized ransomware and remote access trojans (RATs). Early countermeasures within the academic community frequently focused on differential analysis, but pivotal recent works heavily emphasize that attackers have evolved. Salem et al. \cite{b4} highlighted definitively that attackers have heavily adapted by actively crowdsourcing and cataloging the hypervisor traces, memory layout footprints, and instrumentation signatures of even the most heavily cloaked, modern sandboxing solutions. 

To actively effectively counter this intelligence sharing, academic researchers have steadily prioritized increasingly transparent analysis techniques. Pi et al. \cite{b7} demonstrated highly advanced deep learning methodologies designed to dynamically detect environment-aware malware behaviorally without ever needing to trigger deep, highly detectable system software hooks.

\subsection{The `Wear-and-Tear' Check Dilemma}
However, despite enormous technological improvements in root hypervisor transparency and kernel-level un-hooking, simple environmental and user-activity heuristic checks remain a massive, critical unaddressed blind spot across the industry. Vasenkov and Kholod \cite{b3} argued forcefully in their formal taxonomy of sandbox evasion techniques that analytical environments utterly fail out-of-the-box to simulate or mimic authentic "wear-and-tear" artifacts—the deeply embedded, heavily interconnected traces of long-term, genuine user interaction over an extended temporal window.

Al-Ghafari et al. \cite{b2} empirically validated this hypothesis in their defensive perspective analysis, offering incontrovertible proof that modern evasion malware can reliably identify pristine, out-of-the-box sandboxes with staggering accuracy simply by algorithmically observing the mathematical lack of historical user data, utterly untouched offline nested registries, and completely empty SQLite browsing databases. Ferrand \cite{b5} subsequently detailed the bleeding-edge state of the evasion art, noting specifically that environment fingerprinting systems directly targeting these 'human-element' telemetry artifacts constitute the primary defensive setup hurdle for analysts scaling operations today. 

Furthermore, Zimba et al. \cite{b8} rigorously modeled the distinct evolutionary timeline of these specific techniques, statistically demonstrating that while traditional API code obfuscation complexities have somewhat plateaued, complex anti-sandbox environment logic checks have risen exponentially in global frequency metrics. Kumar et al. \cite{b6} subsequently stressed the urgent, unyielding necessity for hyper-realistic behavioral simulation logic in dynamic analysis pipelines to effectively stem this rising tide. 

Our core framework, ARC, builds directly upon the mathematical and theoretical premises heavily established in these recent, critical surveys. Moving firmly and logically beyond mere static hypervisor cloaking efforts, ARC is focused on robust computational generative theory. By programmatically and systematically generating the complex, interconnected `wear-and-tear' topological artifacts of a strictly lived-in system, ARC systematically and actively addresses the modern, data-centric evasion strategies heavily detailed by researchers \cite{b1, b2}.

\section{ARC System Architecture and Design}

The ARC operating framework is stringently engineered on a deeply decoupled, service-oriented micro-architecture explicitly designed for uncompromising robust extensibility, absolute execution determinism, and deep forensic deployment auditability. The entire generation pipeline logically operates completely offline strictly on a mounted, raw VHD file, negating the heavy computational overhead of live boot modification.

\subsection{Core Design Philosophy and Absolute Determinism}
A central tenet of the ARC system is Absolute Determinism. Forensic reproducibility is mandatory for analytical systems. Let $P$ be a configuration profile state, and $S$ be a global initial cryptographic seed. The generation function $G$ modifying a base VHD state $V_0$ to produce personalized state $V_1$ must satisfy:

\begin{equation}
\forall (P, S), G(V_0, P, S) = V_1 \implies G(V_0, P, S) \equiv V_1
\end{equation}

Given the exact same configuration YAML profile definition and the exact identical initialization cryptographic seeds (determining the core username, base starting timestamp, and hardware stochastic indices), the ARC system logically must produce an identically modified binary Windows 11 image output on successive mathematical runs.

\subsection{Pipeline Orchestration Phase Modeling}
The Orchestrator acts as the central execution cerebellum, managing the directed acyclic graph (DAG) of the execution pipeline. It recursively mathematically maps the explicit dependency graph of all registered subsystem services, strictly ensuring that fundamental, core services reliably execute in the correct hierarchical architectural order. For example, underlying NTFS directory structures and core user registry hives must be fundamentally initialized and structured before the application log event generator can successfully utilize them to simulate temporal crashing. The Orchestrator propagates an immutable logical context containing the dictionary profile variables to all descendant abstract services.

\subsection{Profile Engine: Stochastic Persona Synthesis}
Creating a highly believable digital deception intrinsically requires generating a computationally believable, dense user persona. The Profile Engine resolves heavily modular YAML definitions into concrete mathematical usage bounds. A designated "Software Developer" profile algorithmically yields extremely heavy artifact operational densities within specific file path sub-trees, dense JSON configurations, high-entropy complex command-line PowerShell execution histories, and localized Git repository objects. Conversely, an "Enterprise Office User" definition produces fundamentally different spatial outcomes: extreme densities surrounding \texttt{.docx} and \texttt{.xlsx} artifacts, SharePoint URL cached linkages, and complex \texttt{UserAssist} ROT13 keys invariably linked inherently to Microsoft Office binaries.

\subsection{Identity Generator: Algorithmic Resolution}
The Identity Engine creates a deterministic, internally and externally coherent master identity linking both logical User logic and physical Hardware properties. The sub-system parses huge, localized offline dictionaries to randomly generate phonetically and structurally plausible human names, functional organizational corporate domains, randomized internal RFC-1918 IP address schemes, and exact MAC address OUI (Organizationally Unique Identifier) vendor mappings that seamlessly geometrically align with the chosen, synthesized hardware topological profile.

\subsection{The Spatio-Temporal Consistency Engine}
The Timestamp Service acts as the absolute temporal mathematical backbone for the entirety of the framework. Isolated, naïve artifact script generation chronically fails in operations because typical automation tools rapidly assign timestamps indiscriminately, inherently resulting in structural metadata indicating thousands of files were instantiated absolutely simultaneously, within identical milliseconds—an immediate, catastrophic algorithmic red flag.

ARC effectively distributes distinct temporal transaction events stochastically over a widely configurable timeline parameter (e.g., a massive 90-day simulated history window). 

\textbf{Timestamp Generation Algorithm:}
\begin{verbatim}
def gen_logical_timestamp(
    self, base_dt, evt_type, var_mins
):
    """
    Computes a Poisson-distributed 
    temporal event ensuring logical flow.
    """
    # Exclude weekend distribution 
    if self.is_weekend(base_dt) and \
       self.profile_type == 'enterprise':
        base_dt = self.shift_to_monday(base_dt)
        
    # Restrict to typical 0900-1700 dev bells
    time_offset = numpy.random.normal(
        loc=13.0, scale=3.0
    ) 
    
    event_ts = base_dt + timedelta(
        hours=time_offset
    )
    
    # Enforce MAC chronological logic
    if evt_type == 'modification':
        event_ts += timedelta(
            minutes=abs(
                numpy.random.poisson(var_mins)
            )
        )
        
    return event_ts
\end{verbatim}
If a structured document is logged as "edited," the Timestamp Service mathematical logic securely ensures the file's raw underlying \texttt{LastModified} attribute logically, chronologically succeeds its strictly maintained \texttt{CreationTime}, and that the corresponding \texttt{LastAccessTime} makes irrefutable fundamental chronological sequence sense.

\section{Implementation of Artifact Subsystems}

ARC utilizes massive, heavily specialized Python implementations tightly grouped sequentially into extremely distinct logical spatial domains designed to independently interface via structural abstraction interfaces.

\subsection{Immutable Offline Registry Modifications (HiveWriter)}
The Registry Services module, utilizing deep raw binary parsing frameworks specifically tailored for Windows NT registry structures (e.g., \texttt{python-registry} and open-source hive editors), allows ARC to directly edit binary logical structures nested deep within the \texttt{NTUSER.DAT}, \texttt{SOFTWARE}, and \texttt{SYSTEM} proprietary hives directly upon the disk medium, utterly without attempting to boot the OS. This bypasses extensive virtualization overhead. ARC meticulously maps and algorithmically parses standard ROT13 encoding mechanisms traditionally used within \texttt{UserAssist} execution tracking binaries. It updates MRU (Most Recently Used) jump lists, directly injects hyper-realistic "Installed Program" binary keys, and meticulously stages entirely fake persistence mechanisms mathematically simulating highly benign, standard enterprise third-party updaters (e.g., GoogleUpdateTaskMachine).

\subsection{Spatially Distributed Filesystem Emulation}
The deep Filesystem Services topology functions well beyond merely generating blank, empty text documents recursively. The system intelligently constructs highly varied raw document binary types containing deeply genuine, specific profile-relevant natural language textual content dynamically extracted from static datasets. This multi-layered simulation expressly encompasses heavily seeding specific internet media graphical caches, constructing proprietary Windows thumbnails (\texttt{thumbcache\_*.db} binaries), manipulating structurally hidden absolute \texttt{\$Recycle.Bin} logical pointers, and successfully compiling simulated, locked temporary \texttt{\~{}WRL} Microsoft Word proprietary ownership files.

\subsection{Browser SQLite State Synthesis}
Operating dynamically and interactively with the highly proprietary underlying SQLite structured architecture natively utilized by modern, widespread Chromium-based browser clients. The Browser Service synthesizes heavily interwoven, chronologically highly coherent browsing timeline histories (writing directly to the dense \texttt{History} SQLite structures), cleverly stages hundreds of unique third-party JSON tracker cookies, and mathematically populates structurally deep, sound JSON file network architectures for synchronized Bookmarks and internal Configuration Preferences designed perfectly to mimic and simulate extremely localized, active online human network presence.

\subsection{Event Log Sequencing (EVTX Generation)}
Simulated topological file usage is utterly mathematically useless under advanced analysis if the core underlying Windows Event Viewer binary database continuously displays a sterile total machine timeline operational uptime of mathematical zero. Directly utilizing heavily specialized binary \texttt{.evtx} logic parsers targeting Windows XML Event structures, the system forces injection algorithms of deeply interwoven proprietary system and internal security events into the registry logs. It successfully mathematically models and simulates complex baseline topological physical network DHCP connection handshakes over time, simulates extensive physical User Account Control (UAC) administrative consent prompts occurring asynchronously, mimics logical deep Windows Update monthly service patching cycles computationally, and flawlessly generates matching systemic boot, sleep, and core shutdown telemetry mathematical curves.

\section{Advanced Anti-Fingerprinting Strategies}

Beyond the simple, baseline additive layer of complex generative user telemetry data, the ARC framework simultaneously functions as an extraordinarily aggressive topological scrubbing and alteration tool, profoundly intentionally erasing hardcoded structural internal indicators of a standardized synthetic operational environment explicitly tailored to perfectly defeat aggressive system-level baseline malware configuration checks.

\subsection{Strict Hardware Profile Normalization}
Environment-aware malware architectures routinely statically query specific strings contained deep within the \texttt{HKLM\textbackslash HARDWARE\textbackslash DESCRIPTION\textbackslash System} static registry topology constraints upon execution initialization. A generalized standard Cuckoo or VMware ESXi virtual instance dynamically displays highly recognizable, exceptionally conspicuous embedded compiler strings, such as explicitly \texttt{SystemBiosVersion = VBOX\_} or \texttt{VideoBiosVersion = VirtualBox VBE}. 

ARC bypasses this algorithmically by executing a strict, offline global topological normalization string sweep explicitly prior to sandbox mounting. Directly utilizing a deep, constantly-updated comprehensive internal JSON repository payload composed of aggressively verified physical bare-metal hardware enterprise topologies, ARC strictly uniformally statically overwrites deeply-nested, underlying VM-specific motherboard architecture, OEM BIOS string pointers, physical CPU nomenclature parameters, and specific proprietary peripheral strings. It simultaneously dynamically targets and injects correlating matching computational randomly-simulated physical IDE disk drive hardcoded structural volume serial numbers, thoroughly overriding the highly recognizable standardized simplistic virtual abstraction identifiers automatically natively exposed by the computational host system.

\subsection{Targeted Hypervisor Artifact Differential Scrubbing}
The underlying VM Scrubber execution module natively recursively iterates logarithmically deep across the base hierarchical global structural system offline hive, continuously identifying and executing completely targeted systemic deletion topology directives. It aggressively logically isolates, tracks, and then ruthlessly securely purges entirely deeply nested registry configuration logic pointers that inherently functionally map natively unequivocally strictly towards legacy structural hypervisor seamless visual integration orchestration scripts, internal standard operating virtual `Guest Additions' execution components, or specific host-platform memory communication daemon interaction services (for example, aggressively terminating execution remnants of \texttt{VBoxMouse.sys}, \texttt{vmtoolsd.exe}, proprietary VirtualBox shared graphic acceleration, and fundamental underlying Hyper-V Integration Components binary footprint traces).

\subsection{Latent Behavioral Simulation and Faking}
A sophisticated evasive malware zero-day sample may completely choose dynamically to strictly avoid parsing localized static hard-disk files altogether during execution, opting instead logically to heavily functionally probe the live interactive active network and device peripheral driver topological state of the highly complex nested operating system localized core configuration. 

To counter this highly advanced, system-level execution topology behavior, ARC mathematically algorithms systematically natively physically adjust underlying structural deep background operational configurations purely to effectively actively mathematically mimic fundamentally the continuous, active historical presence of fundamentally highly complex baseline standard network physical topological peripheral drivers. It extensively meticulously structurally configures proprietary explicit deep nested static system registry keys mathematically specifically deliberately simply to artificially explicitly effectively dynamically indicate directly that extensive internal structural Bluetooth transmission stacks, complex spatial physical hardware USB audio acceleration processing hardware physical layout sub-layers, and extremely computationally temporally highly complex historical wireless Wi-Fi 802.11 physical enterprise network authentication active physical network profiles have actively dynamically definitely unquestionably historically completely previously been heavily repeatedly consistently massively utilized directly.

\section{Evaluation Methodology}

To quantitatively assert the effectiveness of ARC, we evaluated generated images against a stringent three-tiered validation suite expressly engineered into the project's internal \texttt{evaluation/} directory.

\subsection{Experimental Setup and Baselines}
The controlled experiment utilized a pristine, completely un-activated, sterile clean baseline Windows 11 Enterprise virtual hard disk as the initial standard state control mathematical baseline group topological variable mapping point. 

Three uniquely distinct ARC generation parametric profiles representing divergent entropy topological patterns (the "Software Developer", "Enterprise Office User", and "Personal Home Configuration" datasets) were rigorously automatically executed over the pure baseline identical control disk structural logic systems utilizing a parameterized exactly 90-day execution sequential timeframe.

\subsection{Evaluation Metrics Definition}
We utilized three foundational vector parameters:
\begin{enumerate}
\item \textbf{Geometric Artifact Density Function:} $D(T, A)$ measuring absolute row entries in domain $A$ mapping over time $T$.
\item \textbf{Spatio-Temporal Coherence Rule Index:} A boolean array $C(x, y)$ equating physical mapping constraints vs temporal limitations.
\item \textbf{Detection Resistance Score (DRS):} $DRS = 1.0 - P(\text{Detection})$, where probability maps directly to triggered sandboxing analysis hooks globally.
\end{enumerate}

\section{Results and Analysis}

\subsection{Artifact Density Validation}
The primary mathematical failure execution point of traditional automated enterprise analytical baseline sandboxes directly resides uniquely within explicit statistical base array digital layout geometric data sparsity. We explicitly generated metric logs assessing the baseline logical spatial sequence arrays dynamically executing over Registry, Filesystem, and Chromium Browser operational layers.

\begin{table}[htbp]
\caption{Artifact Density Geometric Measurement Matrix}
\centering
\begin{tabular}{lccc}
\toprule
\textbf{Metric Element Domain} & \textbf{Sterile Base} & \textbf{Typical Real} & \textbf{ARC Baseline} \\
\midrule
Standard Deep User Documents & 0 & $> 1500$ & 2143 \\
Registry \texttt{UserAssist} Base & 3 & $\sim 150$ & 185 \\
SQLite Browser Memory Rows & 0 & $> 5000$ & 8720 \\
Proprietary \texttt{Security.evtx} & 142 & $> 10000$ & 14390 \\
\bottomrule
\end{tabular}
\end{table}

\subsection{Internal Coherence and Spatio-Temporal Consistency}
Across exactly fifty unique rigorously automated independent trial executions, ARC reliably produced an internal baseline execution structural aggregate coherence integrity score of 98.4\%. The system consistently preserved chronological ordering across filesystem creation timestamps and registry last-modified keys, avoiding the timestamp clustering that immediately flags naive simulators.

\subsection{Sandbox Signal Testing (Evasion Resistance)}
We subjected our personalized images directly to a suite of explicitly established sandbox heuristic detection binaries logically derived from open-source analysis evaluation frameworks. The Detection Resistance Score (DRS) rose strictly from a sterile baseline of 0.15 directly to an average of 0.92, successfully defeating the vast majority of standard behavioral environment probes.

\section{Discussion and Limitations}

\subsection{Static Pre-Generation vs Active Engagement}
ARC operates completely offline before execution. While this provides absolute static check resistance, it inherently lacks genuine runtime input simulation (such as active real-time mouse movements or keystrokes during malware detonation). Extremely advanced malware that actively waits for live human interaction metrics over a long duration may still identify the environment.

\subsection{Scalability and OS Limitations}
The deep structural editing of internal binary Windows hive topologies is inherently fragile. Future major revisions to the Windows NT registry format or the proprietary \texttt{.evtx} schema may break ARC's offline manipulation logic, requiring continuous maintenance and mapping of the underlying data structures.

\section{Conclusion}

This research demonstrates that profile-driven VM personalization can fundamentally enhance the effectiveness of dynamic analysis pipelines. By programmatically generating mathematically consistent, dense artifacts across the filesystem, registry, and event logs, ARC effectively simulates realistic human usage. While perfect evasion resistance against active runtime behavioral probing remains an open challenge, this framework drastically raises the difficulty for environment-aware malware, decisively shifting the operational advantage back toward security analysts and automated defense systems.

\section*{References}

\begin{thebibliography}{00}

\bibitem{b1} X. Zhang, J. Xiao, and Y. Guo, ``A Survey of Dynamic Malware Analysis Evasion Techniques and Defenses,'' \textit{IEEE Access}, vol. 11, pp. 23098--23115, 2023.

\bibitem{b2} M. Al-Ghafari, M. Al-Ghafari, and K. El-Khatib, ``Next-Generation Sandbox Evasion: A Taxonomy and Defensive Perspective,'' \textit{International Journal of Information Security}, pp. 1--22, 2024.

\bibitem{b3} A. Vasenkov and I. Kholod, ``A Taxonomy of Sandbox Evasion Techniques,'' in \textit{Proceedings of the 2022 Conference of Russian Young Researchers in Electrical and Electronic Engineering (ElConRus)}, pp. 442--446, 2022.

\bibitem{b4} A. Salem, Y. Boshmaf, and O. Al-Ibrahim, ``Advanced Anti-Analysis Techniques in Modern Malware,'' \textit{IEEE Security \& Privacy}, vol. 21, no. 2, pp. 55--64, 2023.

\bibitem{b5} O. Ferrand, ``State of the Art of Malware Evasion Techniques,'' \textit{Journal of Computer Virology and Hacking Techniques}, vol. 18, no. 1, 2022.

\bibitem{b6} R. Kumar, P. Singh, and K. Sharma, ``Evading Sandbox Analysis: A Comprehensive Survey,'' \textit{Computers \& Security}, vol. 114, 102581, 2022.

\bibitem{b7} L. Pi, Y. Chen, and Z. Zhao, ``A deep learning approach to detect environment-aware malware,'' \textit{Journal of Network and Computer Applications}, vol. 174, 102898, 2021.

\bibitem{b8} A. Zimba, Z. Wang, and H. Chen, ``Modeling the Evolution of Malware Evasion Techniques,'' \textit{Security and Communication Networks}, 2021.

\bibitem{b9} A. Afianian, S. Niksefat, B. Sadeghiyan, and D. Baptiste, ``Malware dynamic analysis evasion techniques: A survey,'' \textit{ACM Computing Surveys (CSUR)}, vol. 52, no. 6, pp. 1--28, 2019.

\bibitem{b10} D. Votipka, A. Sethi, and M. L. Mazurek, ``An Observational Investigation of Reverse Engineers' Processes,'' in \textit{USENIX Security Symposium}, pp. 1875--1892, 2020.

\bibitem{b11} D. C. D’Elia, E. Coppa, V. Palmarini, and L. Cavallaro, ``Tackling Environment-Sensitive Malware with Automated Reasoning,'' in \textit{Proceedings of the IEEE European Symposium on Security and Privacy (EuroS\&P)}, 2020.

\bibitem{b12} R. Sihwail, K. Omar, and K. Z. Ariffin, ``A Survey on Malware Analysis Techniques: Static, Dynamic, Hybrid and Memory Analysis,'' \textit{International Journal on Advanced Science, Engineering and Information Technology}, 2019.

\bibitem{b13} A. Bulazel and B. Yener, ``A survey on automated dynamic malware analysis evasion and counter-evasion: PC, mobile, and web,'' in \textit{Proceedings of the 1st Reversing and Offensive-oriented Trends Symposium}, pp. 1--21, 2017.

\bibitem{b14} N. Miramirkhani, M. P. Applegate, and N. Nikiforakis, ``Spotless sandboxes: Evading malware analysis systems using wear-and-tear artifacts,'' in \textit{2017 IEEE Symposium on Security and Privacy (SP)}, pp. 1009--1024, 2017.

\bibitem{b15} M. Polino, A. Sciancalepore, R. Zito, F. Maggi, and S. Zanero, ``Measuring and Defeating Anti-Analysis Malware,'' in \textit{Annual Computer Security Applications Conference (ACSAC)}, pp. 350--363, 2017.

\bibitem{b16} A. Yokoyama, K. Ishii, R. Makino, K. Yoshioka, S. Shikata, T. Matsumoto, and C. Rossow, ``SandPrint: Fingerprinting malware sandboxes to provide intelligence for sandbox evasion,'' in \textit{International Symposium on Research in Attacks, Intrusions, and Defenses (RAID)}, pp. 165--187, 2016.

\bibitem{b17} D. Kirat and G. Vigna, ``MalGene: Automatic Extraction of Malware Analysis Evasion Signature,'' in \textit{Proceedings of the 22nd ACM SIGSAC Conference on Computer and Communications Security}, pp. 769--780, 2015.

\bibitem{b18} X. Ugarte-Pedrero et al., ``A close look at a daily dataset of malware samples,'' \textit{ACM Transactions on Information and System Security}, 2015.

\end{thebibliography}

\end{document}
